\section{Introduction}
EEG is central to clinical neuroscience---sleep staging, screening for abnormal background
activity, studies of depression and cognitive load---but expert labels are scarce, which
motivates foundation models pretrained once on large unlabeled corpora such as the TUH EEG
Corpus (TUEG) \cite{obeid2016tuh} and transferred to many tasks
\cite{jiang2024labram,wang2025cbramod,csbrain,reve2025}. Their pretraining objectives are
borrowed from other modalities. LaBraM \cite{jiang2024labram} trains a vector-quantized
tokenizer and predicts its discrete codes, as BEiT does for images \cite{bao2022beit}; CBraMod
\cite{wang2025cbramod}, CSBrain \cite{csbrain} and REVE \cite{reve2025} reconstruct the raw
waveform of masked patches, as masked autoencoders reconstruct pixels \cite{he2022mae}. Progress
has come mainly from larger models and more data: REVE pretrains on about 60{,}000 hours.

We present \textbf{SpecBraM}, an EEG foundation model whose pretraining target is a physical
quantity rather than the signal. Spectral power informs many EEG labels: AASM sleep
stages are specified by the amount of delta activity, spindles and alpha
\cite{berry2017aasm}, clinical reports describe ``background slowing'', and quantitative EEG
uses band powers and their ratios. SpecBraM therefore keeps the standard masked-patch
transformer but, for each masked (channel, time) patch, predicts its log power in the five
classical frequency bands, computed from the patch's own spectrum (Fig.~\ref{fig:contrast}). The
target needs no tokenizer and no target network and costs one FFT per patch.

This target makes pretraining efficient, and we motivate a possible mechanism
(Sec.~\ref{sec:theory}). Under squared error, variation remaining after conditioning on the
visible context contributes gradient variance but no expected gradient. In an idealized
stationary-Gaussian Whittle model, marginal Fourier phases are ancillary for spectral
parameters; this does not imply that context cannot predict phase. With exactly flat bands
and observation restricted to those bands, the band-averaged periodogram is sufficient for
the band levels and is their maximum-likelihood estimate. Its zero-floor logarithm has
state-independent conditional sampling variance. These properties motivate a compact spectral
target, but do not prove lower gradient noise or faster optimization for the implemented model.

The experiments establish the efficiency gains. With backbone, data and steps fixed, band power reaches after one to two
pretraining epochs the linear-probe quality that waveform reconstruction reaches after sixteen,
and a band-power model with 4.2M parameters surpasses a waveform model with 51.3M. Pretrained on
2{,}388 hours, SpecBraM is evaluated with an identical strict linear probe on eight datasets
from six task types. Our contributions are:
\begin{itemize}
\item \textbf{SpecBraM}, an EEG foundation model pretrained by masked band-power prediction, a
parameter-free physical target that replaces waveform reconstruction (Sec.~\ref{sec:method}).
\item \textbf{An idealized rationale for the target}: band power summarizes state-associated
spectral levels without requiring phase prediction; risk and gradient decompositions clarify
the role of contextual uncertainty and target noise (Sec.~\ref{sec:theory},
Appendix~\ref{app:theory}).
\item \textbf{Evidence across tasks and in efficiency.} Under a strict linear probe SpecBraM
has higher point estimates than CBraMod, LaBraM and CSBrain on (nearly) all tasks and REVE on
most of them, using less pretraining data than several baselines. Matched target studies show
gains in pretraining steps, model size and low-label sleep probing
(Sec.~\ref{sec:results}).
\end{itemize}

\noindent We do not claim a new paradigm: masked prediction of hand-crafted features is
established in vision \cite{wei2022maskfeat} and speech \cite{hsu2021hubert}, and spectral
information enters the targets of several EEG models (Sec.~\ref{sec:related}).
FAME \cite{yu2026fame} and MANAS-2 \cite{kulkarni2026manas2}, which also use band
log-power targets, are concurrent work.
